\section{Live-Demos und technische Umsetzung}
\label{sec:demo}

\subsection{XSS-Demo}

Die XSS-Demo stellt denselben persistent gespeicherten Wert in einer
anfälligen und einer sicheren Kalenderansicht gegenüber. Der Test
\path{api/src/test/java/com/frostbear/scheduler/demo/XssDemoSetupTest.java}
bereitet eine vorhandene lokale Planungsperiode mit reproduzierbaren
Demo-Daten vor. Die zugehörige Gradle-Aufgabe \path{prepareXssDemo}
und die Property \path{xssDemoPlanningPeriodId} erlauben die gezielte Auswahl.

Die unsichere Ansicht übergibt das Event vollständig an den
Standardrenderer von \texttt{vue-cal}. Die sichere Ansicht
\texttt{ReserveTimeslotSafe.vue} definiert einen eigenen Event-Slot.
Über \texttt{VITE\_XSS\_SAFE\_RESERVATION\_VIEW} und die
Router-Umschaltung kann für die Präsentation zwischen beiden
Darstellungen gewechselt werden.

Der entscheidende Vergleich ist nicht ein unterschiedlicher Payload,
sondern dieselbe gespeicherte Eingabe:

\begin{enumerate}
  \item Unsichere Ansicht öffnen: Der Browser interpretiert den Wert
    als HTML und der lokale Alert erscheint.
  \item Sichere Ansicht öffnen: Der vollständige Payload bleibt
    sichtbar, wird aber nur als Text dargestellt.
  \item Ursache im Code zeigen: Standardrenderer mit HTML-Senke
    gegenüber eigenem Slot mit Vue-Interpolation.
\end{enumerate}

\subsection{CORS-Demo}

Die CORS-Demo verwendet vier lokale Komponenten: Keycloak und
PostgreSQL, die Spring-API auf Port 8079, das reguläre Vue-Frontend
auf Port 8078 sowie die separate Angreifer-Demo aus
\texttt{cors/index.html} auf Port 5173.

Die Seite ruft den öffentlichen Planungsperioden-Export direkt auf
und versucht, die Antwort als Datei zu lesen. Der Vergleich zeigt:

\begin{enumerate}
  \item Ohne passende CORS-Freigabe oder mit fester Allowlist
    blockiert der Browser den Lesezugriff der Origin \texttt{localhost:5173}.
  \item Mit der bewusst unsicheren globalen Freigabe kann die fremde
    Seite die Antwort lesen und den Download auslösen.
\end{enumerate}

\begin{figure}[H]
  \centering
  \begin{minipage}[t]{0.48\textwidth}
    \centering
    \includegraphics[width=\linewidth]{cors-unsecure}
    \caption*{Unsichere globale Freigabe: Antwort ist lesbar.}
  \end{minipage}
  \hfill
  \begin{minipage}[t]{0.48\textwidth}
    \centering
    \includegraphics[width=\linewidth]{cors-secure}
    \caption*{Feste Allowlist: Demo-Origin wird blockiert.}
  \end{minipage}
  \caption{Gegenüberstellung der lokalen CORS-Demo}
  \label{fig:cors-demo}
\end{figure}

\subsection{Aussagekraft und Grenzen}

Die Demos beweisen jeweils einen klar abgegrenzten Mechanismus. Der
XSS-Alert beweist Codeausführung im Scheduler-Origin, aber noch keine
konkrete Datenexfiltration. Der CORS-Dateidownload beweist, dass eine
fremde Origin eine öffentliche Antwort lesen kann, aber noch keinen
Token-Diebstahl.

Insbesondere wird der Bearer-Token des Scheduler-Frontends nicht
automatisch durch \texttt{allowCredentials(true)} an die Demo-Origin
übertragen. Ein vollständiger Account-Takeover würde einen
zusätzlichen Fehler voraussetzen, etwa die nachgewiesene XSS-Klasse
vor ihrem Fix, eine unsichere Token-Antwort oder eine spätere
Cookie-basierte Authentifizierung mit fehlerhaftem CORS.
